\chapter{Ненулевой стационарный режим изолированного стержня}

\section{Постановка задачи и выбор границ}

Пусть стержень занимает отрезок $[0,l]$, где $l>0$, и имеет
теплоизолированную боковую поверхность. Внешние и внутренние источники
тепла отсутствуют. Поэтому в ранее полученном уравнении необходимо
положить $b=0$ и $f=0$:
\[
 u_t=a^2u_{xx},\qquad 0<x<l,\quad t>0,\qquad a>0.
\]
Начальная температура задана условием
\[
 u(x,0)=x.
\]

В первоначальном варианте задания были указаны условия первого рода
при $x=0$ и второго рода при $x=l$. При их однородности получаем
\[
 u(0,t)=0,\qquad u_x(l,t)=0.
\]
Такие условия не дают ненулевого стационарного режима: из
$U''=0$ следует $U=Ax+B$, а условия $U(0)=0$ и $U'(l)=0$
дают $A=B=0$.

Следовательно, выполнить новый пункт с сохранением прежних типов
границ невозможно. Поскольку в этом пункте требуется подобрать
однородные условия, выберем теплоизоляцию обоих торцов:
\[
\boxed{u_x(0,t)=0,\qquad u_x(l,t)=0.}
\]
По закону Фурье нулевая производная температуры означает отсутствие
теплового потока через торец. В обозначениях общей задачи с изображения
этот выбор соответствует
\[
 \alpha=\gamma=0,\qquad\beta=\delta=1,
 \qquad\mu=\nu=0.
\]
Получаем задачу
\[
\boxed{\begin{cases}
 u_t=a^2u_{xx},&0<x<l,\ t>0,\\
 u_x(0,t)=u_x(l,t)=0,&t>0,\\
 u(x,0)=x,&0\le x\le l.
\end{cases}}
\]

\section{Стационарные решения и сохранение тепла}

Стационарное решение не зависит от времени: $u(x,t)=U(x)$.
Для него
\[
 U''=0,\qquad U'(0)=U'(l)=0.
\]
Поэтому
\[
\boxed{U(x)=C,\qquad C\in\mathbb R.}
\]
Таким образом, стационарная краевая задача имеет семейство решений,
включающее ненулевые постоянные. Значение $C$ выбирается начальными данными.

Для произвольной начальной температуры $\phi(x)$ интегрируем уравнение
по длине стержня:
\[
 \frac{d}{dt}\int_0^lu(x,t)\,dx
 =a^2\bigl(u_x(l,t)-u_x(0,t)\bigr)=0.
\]
Следовательно,
\[
 \int_0^lu(x,t)\,dx=\int_0^l\phi(x)\,dx.
\]
При постоянной объёмной теплоёмкости $c$ и площади сечения $S$
величина $cS\int_0^lu\,dx$ выражает тепловую энергию относительно
выбранного температурного отсчёта. В полностью изолированном стержне
она сохраняется.

Если $u(\cdot,t)$ стремится к постоянной $C$, сохранение средней
температуры даёт
\[
 Cl=\int_0^l\phi(x)\,dx,
 \qquad
\boxed{C=\overline\phi=\frac1l\int_0^l\phi(x)\,dx.}
\]
Для $\phi(x)=x$ получаем
\[
\boxed{C=\frac1l\int_0^lx\,dx=\frac l2\ne0.}
\]

\section{Решение методом Фурье}

Рассмотрим задачу на собственные значения
\[
 -X''=\lambda X,\qquad X'(0)=X'(l)=0.
\]
Она имеет нулевую моду
\[
 \lambda_0=0,\qquad X_0(x)=1
\]
и положительные моды
\[
 \lambda_n=\left(\frac{n\pi}{l}\right)^2,\qquad
 X_n(x)=\cos\frac{n\pi x}{l},\qquad n=1,2,\ldots.
\]
Нулевая мода не затухает; все остальные моды затухают экспоненциально.
Для $\phi\in L^2(0,l)$ решение имеет вид
\[
\boxed{u(x,t)=\overline\phi+
 \sum_{n=1}^{\infty}A_n
 e^{-a^2(n\pi/l)^2t}\cos\frac{n\pi x}{l},}
\]
где
\[
 A_n=\frac2l\int_0^l\phi(\xi)\cos\frac{n\pi\xi}{l}\,d\xi.
\]
Ряд задаёт решение при $t>0$, а начальное условие выполняется
в смысле сходимости в $L^2(0,l)$. При достаточной гладкости и
согласовании начальных и граничных данных получается классическое
решение вплоть до начального момента.

Для $\phi(x)=x$ вычисляем коэффициенты:
\[
\begin{aligned}
 A_n&=\frac2l\int_0^lx\cos\frac{n\pi x}{l}\,dx\\
 &=\frac{2l}{n^2\pi^2}\bigl((-1)^n-1\bigr).
\end{aligned}
\]
Следовательно, коэффициенты чётных мод равны нулю, а для нечётных
\[
 A_{2m+1}=-\frac{4l}{(2m+1)^2\pi^2}.
\]
Окончательно,
\[
\boxed{\begin{aligned}
 u(x,t)=\frac l2-\frac{4l}{\pi^2}
 \sum_{m=0}^{\infty}\frac{1}{(2m+1)^2}
 &\exp\left(-\frac{a^2(2m+1)^2\pi^2}{l^2}t\right)\\
 &\times\cos\frac{(2m+1)\pi x}{l}.
\end{aligned}}
\]
При $t\to\infty$ все слагаемые суммы затухают:
\[
\boxed{\lim_{t\to\infty}u(x,t)=\frac l2.}
\]
Более того,
\[
 \sup_{0\le x\le l}\left|u(x,t)-\frac l2\right|
 \le\frac{4l}{\pi^2}e^{-a^2\pi^2t/l^2}
 \sum_{m=0}^{\infty}\frac{1}{(2m+1)^2}
 =\frac l2e^{-a^2\pi^2t/l^2}.
\]
Значит, сходимость к ненулевому стационарному режиму равномерна
по всей длине стержня.

Следует отметить, что начальная функция $\phi(x)=x$ не удовлетворяет
условиям согласования $\phi'(0)=\phi'(l)=0$.
Поэтому не следует требовать непрерывности $u_x$ в угловых точках
$(0,0)$ и $(l,0)$. Это не препятствует существованию решения:
начальная температура восстанавливается при $t\downarrow0$, а
теплоизоляция торцов выполняется для каждого $t>0$.

\section{Необходимое и достаточное условие}

Для выбранных однородных условий Неймана и произвольной
$\phi\in L^2(0,l)$ предельный стационарный режим равен
\[
 U_\infty(x)=\overline\phi.
\]
Поэтому необходимое и достаточное условие его ненулевости имеет вид
\[
\boxed{U_\infty\not\equiv0
 \quad\Longleftrightarrow\quad
 \int_0^l\phi(x)\,dx\ne0.}
\]

Необходимость: если предел равен ненулевой постоянной $C$, то
по сохранению средней температуры
\[
 \int_0^l\phi(x)\,dx=lC\ne0.
\]
Достаточность: если интеграл начальной температуры ненулевой, то
нулевая мода равна $\overline\phi\ne0$, а все остальные моды затухают.
В частности,
\[
 \|u(\cdot,t)-\overline\phi\|_{L^2(0,l)}
 \le e^{-a^2\pi^2t/l^2}\|\phi-\overline\phi\|_{L^2(0,l)}.
\]
Сходимость также равномерна при $t\to\infty$: для любого $t_0>0$
\[
\begin{aligned}
 \sup_x|u(x,t)-\overline\phi|
 &\le e^{-a^2\pi^2(t-t_0)/l^2}
 \left(\sum_{n=1}^{\infty}A_n^2\right)^{1/2}
 \left(\sum_{n=1}^{\infty}
 e^{-2a^2(n\pi/l)^2t_0}\right)^{1/2},\quad t\ge t_0.
\end{aligned}
\]
Правая часть стремится к нулю.

Условие $\phi\not\equiv0$ само по себе недостаточно. Например,
\[
 \phi(x)=x-\frac l2
\]
является ненулевой функцией, но её среднее равно нулю; решение
стремится к нулевому стационарному режиму. Если же
$\phi\ge0$ почти всюду и $\phi$ не равна нулю почти всюду, её
интеграл положителен, и предельный режим ненулевой.

Наконец, существование ненулевых стационарных решений краевой задачи
следует отличать от достижения ненулевого предела конкретным решением.
Постоянные $C\ne0$ допустимы при любых начальных данных, но решение
с нулевым начальным средним стремится именно к $C=0$.
Режим стационарен уже при $t=0$ тогда и только тогда, когда
$\phi$ является постоянной почти всюду; для ненулевого точного
стационарного решения эта постоянная должна быть ненулевой.
